\section{Conclusion}

Nous avons formulé l'exploitation exacte d'une coalition d'opérateurs sur une
fenêtre de décision, avec activations binaires, capacités et coûts
hétérogènes, puis relié son coût optimal à un jeu coopératif d'économies. La
sous-additivité du coût rend ce jeu superadditif, sans garantir son
équilibrage : les seuils de capacité peuvent produire un cœur vide malgré
l'intérêt collectif de la grande coalition. Les conditions suffisantes, les
certificats analytiques et les indicateurs numériques proposés permettent de
distinguer l'efficacité opérationnelle de la stabilité des allocations.

Le régime de très faible trafic illustre une seconde distinction. La capacité
de chaque équipement à servir seul toute la demande garantit un cœur non vide,
mais ne suffit pas à y placer la valeur de Shapley lorsque les coûts fixes sont
hétérogènes. L'appartenance de Shapley au cœur doit donc être testée
explicitement dans les instances hétérogènes ; l'homogénéité des coûts fixes et
variables fournit, dans ce régime, une condition suffisante simple qui
rétablit cette stabilité.

L'évaluation semi-empirique devra enfin quantifier la fréquence et l'intensité
de ces phénomènes sur les sites multi-opérateurs construits à partir des
mesures disponibles. L'extension à l'incertitude du trafic, aux coûts de
transition et aux interactions multisites constitue une suite naturelle de ce
travail.
